\chapter{Orçamento}

Durante a etapa de orçamento do projeto, buscou-se minimizar os custos de forma a garantir a viabilidade financeira para todos os integrantes do grupo. Para isso, foram realizados levantamentos de preços referentes aos componentes de hardware, estrutura e consumo energético, além da matéria-prima necessária para a fabricação do chassi do micromouse e do labirinto de testes.

Houve uma redução significativa nos custos totais devido à disponibilidade prévia de alguns componentes por parte dos membros da equipe, o que nos permitiu direcionar uma parcela maior do orçamento para a aquisição de componentes com maior robustez, contribuindo para a melhoria da qualidade do projeto.

\vspace{0.5cm}

\begin{table}[htpb]
\centering
\caption{Orçamento geral.}
\begin{tabular}{|l|c|c|}
\hline
\textbf{Geral}          & \textbf{Preço previsto (R\$)} & \textbf{Preço realizado (R\$)} \\ \hline
Matéria-prima           & 84,53                         & 42,50                          \\ \hline
Componentes eletrônicos & 695,67                        & 486,04                         \\ \hline
Montagem                & 330,84                         & 317,34                          \\ \hline
\end{tabular}
\end{table}

\vspace{0.5cm}

Na tabela de orçamento por aquisição, apresentada abaixo, os componentes cujo custo real é igual a zero correspondem aos itens disponibilizados pelos integrantes do grupo.

\vspace{0.5cm}

\begin{landscape}
\begin{table}[htpb]
\centering
\caption{Orçamento por aquisição.}
\begin{tabular}{|l|l|c|c|c|}
\hline
\textbf{Item ou serviço} & \textbf{Tipo} & \textbf{Preço previsto (R\$)} & \textbf{Preço realizado (R\$)} & \textbf{Responsável} \\ \hline
Motor                    & Componente eletrônico & 151,80 & 0,00   & Gabriela \\ \hline
Sensor ToF               & Componente eletrônico & 42,45  & 61,19  & Gabriela \\ \hline
Ponte H                  & Componente eletrônico & 50,15  & 36,00   & Gabriela e Isaque \\ \hline
Giroscópio               & Componente eletrônico & 16,05  & 28,05  & Gabriela \\ \hline
Regulador de tensão      & Componente eletrônico & 7,12   & 0,00   & Gabriela \\ \hline
ESP32-S3                 & Componente eletrônico & 297,80 & 297,80 & Gabriela e Isaque \\ \hline
Rodas                    & Montagem              & 99,80  & 106,30  & Gabriela \\ \hline
MDF                      & Matéria-prima         & 84,53  & 42,50  & Isaque   \\ \hline
Bateria                  & Componente eletrônico & 67,80  & 0,00   & Gabriela \\ \hline
Chassi                   & Montagem              & 98,00  & 78,00  & Gabriela e Davi \\ \hline
Tintas                   & Montagem              & 17,64  & 17,64  & Gabriela \\ \hline
Roda boba                & Montagem              & 17,50  & 17,50  & Gabriela \\ \hline
Jumpers e barra de pinos & Componente eletrônico & 35,00  & 35,00  & João \\ \hline
Placa perfurada & Componente eletrônico & 11,40  & 11,40  & Davi \\ \hline
Insertos metálicos & Montagem & 22,90  & 22,90  & Gabriela \\ \hline
Barra de pinos & Componente eletrônico & 15,94  & 15,94  & Davi \\ \hline
Resistor e diodo & Componente eletrônico & 0,66  & 0,66  & Davi \\ \hline
Cabo de rede & Montagem & 75,00  & 75,00  & Ludmila \\ \hline
\end{tabular}
\end{table}

\begin{table}[htpb]
\begin{center}
\caption{Justificativas para itens ou serviços trocados após a aquisição.}
\begin{tabular}{|c|c|} \hline
\textbf{Item ou serviço} & \textbf{Justificativa}  \\ \hline 
   Chassi &
   Houve erros de dimensionamento nas primeiras versões. \\ \hline
   ESP32-S3 &
   Algumas ESPs foram queimadas. \\ \hline
\end{tabular}
\end{center}
\end{table}

\end{landscape}
